\section{Introduction}
\label{sec:intro}

Many scientific comparisons ask whether two groups of units generate the same shape. In ecology, a site may contribute a cloud of species trait coordinates and the question is whether two habitat classes support the same morphological disparity; in biomedicine, a patient may contribute a cloud of cell locations and the question is whether disease alters tissue architecture. In each case the unit is one point cloud, summarized by a diagram $D = \Dgm_d(F(Y))$, with a functional summary such as the power-weighted silhouette as the working representation \cite{kim_lee_tce_2026,chazal_socg_2014,bubenik_jmlr_2015,berry_jact_2020}.

Units are rarely exchangeable across groups. They carry covariates $X$ such as site, age, batch, or temperature, and those covariates often determine group membership, because they influence who is sampled, treated, or diagnosed. The same covariates shape the topology of the cloud, because the geometry of a sample depends on where and how it was collected. Let $A$ denote the group label and $e(x) = \Prob(A=1 \mid X=x)$ the propensity. When both $e$ and the conditional law of the diagram given $X$ depend on $X$, the groups differ in covariate composition, and a comparison that ignores $X$ compares mixtures rather than like with like. A morphometric study may find more loops in one basin because sampled elevations differ on average, not because basin membership changes shape, and an imaging study may find different persistence profiles because scanning sites differ in resolution and in the fraction of patients scanned.

Almost all existing topological two-sample tests compare the observational conditional laws, $H_0^{\mathrm{cond}}: \Law(D \mid A=1) = \Law(D \mid A=0)$. Robinson and Turner test equality of within-group bottleneck distance distributions by permutation \cite{robinson_turner_2017}; Kwitt et al. use a persistence scale-space kernel with a kernel maximum mean discrepancy test \cite{kwitt_neurips_2015,gretton_2012}; Han et al. compare weighted persistence intensity functions \cite{han_kim_kim_2026}; STRAND maps persistence features to survival functions \cite{murris_strand_2026}; Moon and Lazar test pixel means of persistence images \cite{moon_lazar_2023}; Dubey and Muller use a Fr\'echet analysis of variance \cite{dubey_muller_2019}; and Krebs and Rademacher target relevant differences \cite{krebs_rademacher_2024}. The difficulty is that under covariate shift this null is a statement about covariate composition, so a rejection can be caused by imbalance rather than a topological difference, and a non-rejection can hide a difference within covariate strata.

We prove that both failures occur in explicit, non-degenerate families. First, a false-positive result: a continuous process in which a covariate drives the group label through a logistic propensity and the loop count through a monotone topology knob, the potential-outcome laws coincide exactly, so the causal nulls hold, while the conditional laws separate. At the strongest imbalance the marginal loop-count gap has magnitude 0.373 (exact quadrature 0.3727 and simulation -0.3738 with Monte Carlo standard error 0.0013). The six tests above then reject with probability tending to one; in a sweep of 1000 replications at 400 units, five of the six reach rejection rates between 0.98 and 1.00 at the strongest imbalance while the slowest reaches 0.651. Second, a masking result: a continuous process with strictly positive propensity in which the marginal conditional laws are exactly equal, so every valid level-alpha test of $H_0^{\mathrm{cond}}$ has power at most alpha and no level-$\alpha$ test of that null can be consistent against this covariate-adjusted alternative, while the covariate-standardized silhouette contrast is $\sup_t \abs{\psi_d(t)} = 0.4934$ (batch standard error $2\times10^{-4}$) and the distribution-level contrast is nonzero. The effect is real but invisible to any procedure that compares pooled arm laws, while a covariate-adjusted analysis detects it with power near one. These failures were demonstrated with a deliberately uncalibrated prototype whose size, between 0.005 and 0.015, fell outside the pre-registered operating band; the corresponding gate is recorded as inconclusive rather than certified, and the calibrated cross-fitted test developed here supersedes that prototype. The leakage conclusions rest on the constructions and competitor behavior, not on the prototype.

We recast the comparison as a covariate-standardized topological contrast in a potential-outcomes framework. Let $Y^a$ denote the cloud a unit would have under group $a$, let $D^a = \Dgm_d(F(Y^a))$, and assume causal consistency, arm-specific conditional exchangeability $Y^a \perp A \mid X$, and positivity. The g-formula then identifies the interventional diagram law as a covariate mixture of the observed conditional laws, and with it the standardized contrasts used here, while the marginal topological effect is not identified from conditional topological ignorability alone unless weak exchangeability holds and the representation commutes with mixing \cite{kim_lee_tce_2026,souto_diamantis_tcda_2026,saki_faghihi_2026}. We target two nulls. The outcome-level null $H_0^{\mathrm{out}}$ states that the mean potential silhouette curve is equal in the two arms. The distribution-level null $H_0^{\mathrm{dist,grid}}$ states that the expected persistence measures of the two interventional laws agree after projection onto one fixed shared grid. The outcome-level test is a cross-fitted augmented inverse probability weighting procedure \cite{chernozhukov_dml_2018,new_robins_2018,kennedy_drlearner_2023} whose statistic is a maximum over homology degrees and scales, calibrated by Gaussian multipliers \cite{cck_2013}. The distribution-level test is a studentized maximum over grid coordinates of the projected expected measure, with shared multipliers and Phipson and Smyth correction \cite{phipson_smyth_2010}. The nulls are representation-relative: for the normalized silhouette paired with the expected measure they are logically independent, while a characteristic kernel mean embedding makes the kernel distribution-level null strictly stronger \cite{kwitt_neurips_2015}. Every comparison below names the pairing it belongs to.

Five contributions follow. First, we give a target-relative comparison protocol. It states which quantities may be compared, shows that rejection rates computed under different nulls are not efficiency statements, and requires any claim of superiority to be inside a pre-declared target-matched class. Alongside it sit the false-positive and masking theorems described above: explicit continuous data-generating processes, exact mixture algebra for the masking construction, and consistency lemmas for the six observational tests under the constructed alternatives. Together the protocol and the leakage results explain when an unadjusted rejection is informative and when it is an artefact of imbalance.

Second, we develop the cross-fitted doubly robust outcome-level test and document its calibration scope. The test estimates the mean potential silhouette curve with AIPW scores computed across folds, takes a supremum over scales and a maximum over homology degrees, and calibrates with a Gaussian multiplier process whose multipliers are shared across degrees to preserve dependence. We prove the corresponding local limit and efficiency properties and report where finite-sample calibration departs from the asymptotic prescription. Frozen-label permutation is retained as a diagnostic only, because with estimated nuisances it is not a finite-sample exact test; the production calibration is the studentized multiplier construction.

Third, we build the distribution-level test on the expected persistence measure and state its qualification. The measure is computed on one frozen shared grid, the statistic is a studentized maximum over grid coordinates, and coordinates whose variance falls below a floor are dropped and counted. We prove fixed-grid asymptotic validity. The weak-null size lies in the pre-registered band for every gating cell with at least 200 units except the strongest confounding setting at 200 units, where it reads 0.096 (standard error 0.013); the same cell reads 0.066 at 500 units. We therefore recommend at least 500 units and claim no uniform small-sample control. Equality of the unprojected measures is not certified by a finite grid, and the test targets the projection.

Fourth, we analyze local power, efficiency, and asymptotic relative efficiency among target-matched tests. We derive a local limit and compare it with a pre-registered Gaussian-shift prediction. That screen returns confirmed=false: at 500 units the observed rejection rate is 0.730 against a predicted 0.628, a gap of 0.102 with Monte Carlo standard error 0.026, and a full-fleet check reports passed=false, with 9 of 51 scored rows outside the 0.06 tolerance. A finite-n correction, combining out-of-span truncation with an estimated propensity mean ratio, moves the headline gap to -0.053, but its mean-ratio rate is not proved, so we report both numbers and the criterion behind each. At fixed grid dimension we derive the efficient influence function and semiparametric efficiency bound for both targets, at the estimator level rather than for testing power; four of seven numerical criteria show finite-sample discrepancies, which we report. Among target-matched tests we prove asymptotic relative efficiency one in the non-prognostic design and variance dominance under stated conditions, while recording the finite-sample caveat that in the exact-design record the adjusted test's power sits 0.14 to 0.19 below the unadjusted test at 200 and 500 units; the asymptotic identity does not translate into finite-sample power equality. Power monotonicity of the studentized maximum under positive semidefinite covariance order remains a conjecture supported by numerical probes, and the minimax comparison is limited to a Le Cam bound and a rate slope; no dominance over existing intensity-based tests is claimed.

Fifth, we report a replicated-cloud bake-off with named wins and losses. The adjusted tests hold size under randomization and under strong confounding at 200 units, where the six observational tests reject between 0.43 and 0.95. The adjusted outcome test is late at small samples and under confounding: on a mean shift under randomization at 50 units it holds 0.110 against 0.800 for Robinson-Turner and 1.000 for STRAND, maximum mean discrepancy, and Han et al., and on the full-stage mean and strong confounding cell at 100 units it holds 0.332 against 0.996 for STRAND and 0.992 for Han et al. Distribution-level tests reach near-full power on deterministic separation and dispersion alternatives, and the studentized maximum detects location shifts even at 50 units; the kernel maximum mean discrepancy is blind to a rare-feature alternative, and Han et al. win on faint intensity signals. Competitor rejection rates under confounding mix signal with size distortion, so no cross-null ranking is drawn. The bake-off does not cover a 500-unit loops-cloud operating point; the only 500-unit evidence is a deterministic pilot.

What this paper does not claim. The evidence is synthetic: there is no real-data analysis, and every reported rate comes from simulation with its replication count and standard error. The distribution-level conclusion concerns one frozen finite grid, not the unprojected measure. We do not claim blanket dominance; the adjusted tests lose to existing tests at small samples and under randomization. We do not claim power optimality: monotonicity is a conjecture, the minimax constant is conjectured while only the Le Cam bound and rate slope are proved, and the efficiency statements are estimator-level at fixed grid dimension. We do not claim finite-sample exactness for the multiplicity procedure, whose studentized comparator is an adaptation carrying one conservative finite-sample miss, and we do not claim that the pre-registered local-power criterion is confirmed or that the asymptotic relative efficiency identity is numerically realized at the sample sizes studied.

The paper proceeds as follows. The next section fixes the estimands and identification conditions and develops the two test statistics; the theory sections prove the target-relative protocol and leakage theorems, then the local-limit, efficiency, finite-grid, and minimax results; the experiments section reports the bake-off; and the final section discusses limitations. Proofs and supplementary tables are collected in the appendices.

\section{Related Work}
\label{sec:related}

Topological two-sample testing is the closest strand. Robinson and Turner \cite{robinson_turner_2017} introduced a permutation test based on within-group bottleneck distances. Kwitt et al. \cite{kwitt_neurips_2015} proved universality of the exponentiated persistence scale-space kernel, turning the problem into a kernel maximum mean discrepancy test \cite{gretton_2012}. Han et al. \cite{han_kim_kim_2026} compare weighted persistence intensity functions with a permutation U-statistic. STRAND \cite{murris_strand_2026} maps persistence features to survival functions. Moon and Lazar \cite{moon_lazar_2023} test pixel means of persistence images. Dubey and Muller \cite{dubey_muller_2019} propose a Fr\'echet analysis of variance for random objects, although diagram space is nonnegatively curved and Fr\'echet means are not unique, so the curvature conditions of their central limit theorem require separate justification \cite{mileyko_2011,turner_2014,che_jact_2024}. Krebs and Rademacher \cite{krebs_rademacher_2024} target relevant differences rather than exact equality. All of these test the observational conditional null. They are appropriate when units are exchangeable or when that law is the target, and several are consistent against fixed alternatives; our concern is the target rather than the procedures' validity under their own null, and the targets diverge when covariates drive both membership and topology.

Causal inference and doubly robust estimation provide the estimation layer. Kim and Lee \cite{kim_lee_tce_2026} introduced topological causal effects with the silhouette representation and identified them under conditional exchangeability and positivity. Souto and Diamantis \cite{souto_diamantis_tcda_2026} developed g-formula identification of interventional diagram laws and representation-relative distribution-level targets. Saki and Faghihi \cite{saki_faghihi_2026} prove the complementary negative result that the marginal topological effect is not identified from conditional topological ignorability alone when the summary is non-injective. We build on that identification layer and concentrate on testing. The estimation machinery is the double machine learning and cross-fitting tradition of Chernozhukov et al. \cite{chernozhukov_dml_2018} and Newey and Robins \cite{new_robins_2018}, the doubly robust targeted learning of Kennedy \cite{kennedy_drlearner_2023,kennedy_review_2022}, and the role of the estimated propensity score in efficient estimation \cite{hirano_imbens_ridder_2003,hahn_1998}. The contribution here is to make the target explicit, derive the influence functions for the resulting topological functionals, and separate legitimate comparisons from illegitimate ones.

Functional and silhouette inference supplies the population objects and their limit theory. Chazal et al. \cite{chazal_socg_2014} proved stochastic convergence of persistence landscapes and silhouettes, Bubenik \cite{bubenik_jmlr_2015} developed their statistical theory including a central limit theorem, and Berry et al. \cite{berry_jact_2020} survey functional summaries of persistence diagrams. Divol and Chazal \cite{divol_chazal_2020} study the density of expected persistence diagrams, and Divol and Lacombe \cite{divol_lacombe_2021} develop the geometry of persistence measure space. These results underpin the functional averages and expected measures that our tests standardize; we place them inside a causal contrast and calibrate against a weak null.

Multiple testing over topological features is the fourth strand. Vejdemo-Johansson and Mukherjee \cite{vejdemo_mukherjee_2022} extended multiple hypothesis testing to persistent homology with family-wise error rate control over families of homology features. We adapt that line to a family of grid coordinates of the projected expected persistence measure and, in the outcome-level test, of homology degrees and scales. Our comparator is a studentized maximum with shared multipliers and pooled studentization rather than the original procedure. The adaptation controls the family-wise error rate in the settings we study, with one conservative finite-sample miss that falsifies exact finite-sample equality of size rather than asymptotic control; no finite-sample exactness is claimed.
